% GENERATED FILE. Edit canonical lessons, then run npm run generate:books.
% canonical-source-hash: fnv1a64:84933f7fdf3f8060
% canonical-lessons: GU-C314-hadtal, GU-C314-morco, GU-C314-sutra, GU-C314-pravacan, GU-C314-gunjvu

\chapter{The Strike}
\label{ch:gu-the-strike}
\hlchaptermodality{314}

\begin{chapteropening}
\textbf{By the end of this chapter:} \emph{I can say why people went on strike and marched.}

Complete the last lesson of chapter 314: I can say why people went on strike and marched.
\end{chapteropening}

\section[\texorpdfstring{\gu{હડતાળ} (haḍtāḷ)}{હડતાળ (haḍtāḷ)}]{\gu{હડતાળ} (haḍtāḷ) — a strike (stopping work)}

\label{lesson:GU-C314-hadtal}



\begin{quote}
Before the new one: say the Gujarati for information, then the Gujarati for word (of something).
\end{quote}

\subsection*{You'll want to know: \gu{હડતાળ}}
\textbf{\gu{હડતાળ}} — \emph{haḍtāḷ} — \textquotedblleft{}a strike (stopping work)\textquotedblright{}.

\begin{quote}
\gu{હડતાળથી} \gu{લોકોને} \gu{ખૂબ} \gu{દુઃખ} \gu{થયું}.
\end{quote}

\begin{quote}
\emph{haḍtāḷthī lokone khūb duḥkh thayũ.}
\end{quote}

\begin{quote}
The strike caused people a lot of hardship.
\end{quote}

A feminine noun.

\begin{cousinweb}
\textbf{\gu{હડતાળ}} (\emph{haḍtāḷ}) is probably built from \emph{haṭ}, a market, from Sanskrit \emph{haṭṭa}, where CDIAL lists the Hindi cousin \emph{haṛtāl}, the closing of shops in protest (CDIAL 13944), and \textbf{\gu{તાળું}}, a lock (CDIAL 5749): the market locked up. Work stops: \textbf{\gu{ડ્રાઈવરોની} \gu{હડતાળ} \gu{છે}}, the drivers are on strike.
\end{cousinweb}

\begin{grammarlens}[title={What you've built}]
One more word for telling what happened.
\end{grammarlens}

\subsection*{Your turn}
\begin{itemize}
\raggedright
  \item \emph{Say it:} \emph{haḍtāḷ}
  \item \emph{Say it:} \emph{haḍtāḷ}, once more
  \item \emph{Say it:} \emph{haḍtāḷthī lokone khūb duḥkh thayũ.} — The strike caused people a lot of hardship.
  \item \emph{Recall:} say the Gujarati for information, then the Gujarati for word (of something), then say \emph{haḍtāḷ} again
\end{itemize}

\begin{tcolorbox}[breakable,colback=teal!4,colframe=teal!35!black,title={Before you move on}]
What does \gu{હડતાળ} mean? (\textquotedblleft{}A strike (stopping work)\textquotedblright{}.) Say it once more.
\end{tcolorbox}
\section[\texorpdfstring{\gu{મોરચો} (morco)}{મોરચો (morco)}]{\gu{મોરચો} (morco) — a protest march}

\label{lesson:GU-C314-morco}



\begin{quote}
Before the new one: say the Gujarati for word (of something), then the Gujarati for a strike.
\end{quote}

\subsection*{You'll want to know: \gu{મોરચો}}
\textbf{\gu{મોરચો}} — \emph{morco} — \textquotedblleft{}a protest march\textquotedblright{}.

\begin{quote}
\gu{લોકોએ} \gu{શહેરમાં} \gu{મોટો} \gu{મોરચો} \gu{કાઢ્યો}.
\end{quote}

\begin{quote}
\emph{lokoe śahermā̃ moṭo morco kāḍhyo.}
\end{quote}

\begin{quote}
People held a big protest march in the city.
\end{quote}

A masculine noun.

\begin{cousinweb}
\textbf{\gu{મોરચો}} (\emph{morco}) is usually traced to Persian \emph{mōrča}, \textquotedblleft{}an entrenchment\textquotedblright{}, a soldiers' word: in Gujarati it is a march of protest, which people \textbf{\gu{કાઢે} \gu{છે}}, bring out into the street.
\end{cousinweb}

\begin{grammarlens}[title={What you've built}]
One more word for telling what happened.
\end{grammarlens}

\subsection*{Your turn}
\begin{itemize}
\raggedright
  \item \emph{Say it:} \emph{morco}
  \item \emph{Say it:} \emph{morco}, once more
  \item \emph{Say it:} \emph{lokoe śahermā̃ moṭo morco kāḍhyo.} — People held a big protest march in the city.
  \item \emph{Recall:} say the Gujarati for word (of something), then the Gujarati for a strike, then say \emph{morco} again
\end{itemize}

\begin{tcolorbox}[breakable,colback=teal!4,colframe=teal!35!black,title={Before you move on}]
What does \gu{મોરચો} mean? (\textquotedblleft{}A protest march\textquotedblright{}.) Say it once more.
\end{tcolorbox}
\section[\texorpdfstring{\gu{સૂત્ર} (sūtra)}{સૂત્ર (sūtra)}]{\gu{સૂત્ર} (sūtra) — a slogan}

\label{lesson:GU-C314-sutra}



\begin{quote}
Before the new one: say the Gujarati for a strike, then the Gujarati for a protest march.
\end{quote}

\subsection*{You'll want to know: \gu{સૂત્ર}}
\textbf{\gu{સૂત્ર}} — \emph{sūtra} — \textquotedblleft{}a slogan\textquotedblright{}.

\begin{quote}
\gu{લોકો} \gu{સૂત્રો} \gu{બોલતા} \gu{હતા}.
\end{quote}

\begin{quote}
\emph{loko sūtro boltā hatā.}
\end{quote}

\begin{quote}
People were shouting slogans.
\end{quote}

A neuter noun.

\begin{cousinweb}
\textbf{\gu{સૂત્ર}} (\emph{sūtra}) is a tatsama, Sanskrit \emph{sūtra}, \textquotedblleft{}a thread\textquotedblright{}, and then a short rule held in a line (CDIAL 13561). In a march it is the short line everyone shouts together.
\end{cousinweb}

\begin{grammarlens}[title={What you've built}]
One more word for telling what happened.
\end{grammarlens}

\subsection*{Your turn}
\begin{itemize}
\raggedright
  \item \emph{Say it:} \emph{sūtra}
  \item \emph{Say it:} \emph{sūtra}, once more
  \item \emph{Say it:} \emph{loko sūtro boltā hatā.} — People were shouting slogans.
  \item \emph{Recall:} say the Gujarati for a strike, then the Gujarati for a protest march, then say \emph{sūtra} again
\end{itemize}

\begin{tcolorbox}[breakable,colback=teal!4,colframe=teal!35!black,title={Before you move on}]
What does \gu{સૂત્ર} mean? (\textquotedblleft{}A slogan\textquotedblright{}.) Say it once more.
\end{tcolorbox}
\section[\texorpdfstring{\gu{પ્રવચન} (pravacan)}{પ્રવચન (pravacan)}]{\gu{પ્રવચન} (pravacan) — a talk; a speech}

\label{lesson:GU-C314-pravacan}



\begin{quote}
Before the new one: say the Gujarati for a protest march, then the Gujarati for a slogan.
\end{quote}

\subsection*{You'll want to know: \gu{પ્રવચન}}
\textbf{\gu{પ્રવચન}} — \emph{pravacan} — \textquotedblleft{}a talk; a speech\textquotedblright{}.

\begin{quote}
\gu{સાંજે} \gu{ગુરુએ} \gu{ચોકમાં} \gu{પ્રવચન} \gu{આપ્યું}.
\end{quote}

\begin{quote}
\emph{sā̃je gurue cokmā̃ pravacan āpyũ.}
\end{quote}

\begin{quote}
In the evening the teacher gave a speech in the square.
\end{quote}

A neuter noun.

\begin{cousinweb}
\textbf{\gu{પ્રવચન}} (\emph{pravacan}) is a tatsama, Sanskrit \emph{pravacana}, \textquotedblleft{}speaking, teaching, exposition\textquotedblright{} (Monier-Williams). A speech is given: \textbf{\gu{પ્રવચન} \gu{આપ્યું}}.
\end{cousinweb}

\begin{grammarlens}[title={What you've built}]
One more word for telling what happened.
\end{grammarlens}

\subsection*{Your turn}
\begin{itemize}
\raggedright
  \item \emph{Say it:} \emph{pravacan}
  \item \emph{Say it:} \emph{pravacan}, once more
  \item \emph{Say it:} \emph{sā̃je gurue cokmā̃ pravacan āpyũ.} — In the evening the teacher gave a speech in the square.
  \item \emph{Recall:} say the Gujarati for a protest march, then the Gujarati for a slogan, then say \emph{pravacan} again
\end{itemize}

\begin{tcolorbox}[breakable,colback=teal!4,colframe=teal!35!black,title={Before you move on}]
What does \gu{પ્રવચન} mean? (\textquotedblleft{}A talk; a speech\textquotedblright{}.) Say it once more.
\end{tcolorbox}
\section[\texorpdfstring{\gu{ગુંજવું} (gunjvũ)}{ગુંજવું (gunjvũ)}]{\gu{ગુંજવું} (gunjvũ) — to resound}

\label{lesson:GU-C314-gunjvu}



\begin{quote}
Before the new one: say the Gujarati for a slogan, then the Gujarati for a speech.
\end{quote}

\subsection*{You'll want to know: \gu{ગુંજવું}}
\textbf{\gu{ગુંજવું}} — \emph{gunjvũ} — \textquotedblleft{}to resound\textquotedblright{}.

\begin{quote}
\gu{આખો} \gu{ચોક} \gu{સૂત્રોથી} \gu{ગુંજતો} \gu{હતો}.
\end{quote}

\begin{quote}
\emph{ākho cok sūtrothī gunjto hato.}
\end{quote}

\begin{quote}
The whole square was ringing with slogans.
\end{quote}

Past (it …): \gu{તે} \gu{ચોક} \gu{ગુંજ્યો}, \emph{te cok gunjyo}.

\begin{cousinweb}
\textbf{\gu{ગુંજવું}} (\emph{gunjvũ}) goes back to Sanskrit \emph{guñja}, \textquotedblleft{}humming\textquotedblright{}, a sound word; CDIAL lists the Hindi cousin \emph{gū̃j}, an echo, but not the Gujarati verb (CDIAL 4173). A place \textbf{\gu{ગુંજે} \gu{છે}}, rings, with voices or music.
\end{cousinweb}

\begin{grammarlens}[title={What you've built}]
One more word for telling what happened.
\end{grammarlens}

\subsection*{Your turn}
\begin{itemize}
\raggedright
  \item \emph{Say it:} \emph{gunjvũ}
  \item \emph{Say it:} \emph{gunjvũ}, once more
  \item \emph{Say it:} \emph{ākho cok sūtrothī gunjto hato.} — The whole square was ringing with slogans.
  \item \emph{Recall:} say the Gujarati for a slogan, then the Gujarati for a speech, then say \emph{gunjvũ} again
\end{itemize}

\begin{tcolorbox}[breakable,colback=teal!4,colframe=teal!35!black,title={Before you move on}]
What does \gu{ગુંજવું} mean? (\textquotedblleft{}To resound\textquotedblright{}.) Say it once more.
\end{tcolorbox}
